\documentclass{article}

\usepackage{amsmath,amssymb}
\usepackage{xurl}
\usepackage[hidelinks]{hyperref}
\setlength{\emergencystretch}{2em}

% Shared commands for notes in docs/paper-gaps/.

\DeclareUnicodeCharacter{2080}{\ifmmode{}_{0}\else\textsubscript{0}\fi}
\DeclareUnicodeCharacter{2081}{\ifmmode{}_{1}\else\textsubscript{1}\fi}
\DeclareUnicodeCharacter{2082}{\ifmmode{}_{2}\else\textsubscript{2}\fi}
\DeclareUnicodeCharacter{2099}{\ifmmode{}_{n}\else\textsubscript{n}\fi}

\newcommand{\Lean}{\textsc{Lean}}
\ExplSyntaxOn
\NewDocumentCommand{\leanid}{m}{
  \ifmmode
    \textsf{\detokenize{#1}}
  \else
    \group_begin:
    \urlstyle{sf}
    \tl_set:Nn \l_tmpa_tl {#1}
    \tl_replace_all:Nnn \l_tmpa_tl {\_} {_}
    \exp_args:NV \path \l_tmpa_tl
    \group_end:
  \fi
}
\ExplSyntaxOff
\providecommand{\tr}{\operatorname{tr}}
\newcommand{\tnleanIssueBase}{https://github.com/LionSR/TNLean/issues/}
\newcommand{\tnleanPrBase}{https://github.com/LionSR/TNLean/pull/}
\newcommand{\tnleanDocBase}{https://sirui-lu.com/TNLean/docs/}
\newcommand{\ghissue}[1]{\href{\tnleanIssueBase#1}{GitHub issue \##1}}
\newcommand{\ghpr}[1]{\href{\tnleanPrBase#1}{PR \##1}}
\newcommand{\doclink}[1]{\href{\tnleanDocBase#1.html}{\path{#1}}}

% Dirac notation and the identity operator, as in blueprint/src/macros/common.tex.
% \providecommand keeps note-local definitions authoritative.
\usepackage{dsfont}
\providecommand{\ket}[1]{|#1\rangle}
\providecommand{\bra}[1]{\langle #1|}
\providecommand{\braket}[2]{\langle #1 | #2 \rangle}
\providecommand{\ketbra}[2]{|#1\rangle\!\langle #2|}
\providecommand{\Id}{\mathds{1}}
\providecommand{\one}{\mathds{1}}
\providecommand{\C}{\mathbb{C}}
\providecommand{\MN}[1]{M_{#1}(\C)}
\providecommand{\MD}{M_D(\C)}

% External referencing for paper sources.  The build helper writes sanitized
% auxiliary files under build/paper-aux/ and excludes duplicated source labels.
\usepackage{xr}

% Import a generated paper auxiliary file with a caller-supplied label prefix.
% The candidate paths support builds from docs/paper-gaps/, the repository root,
% and build/paper-aux/ itself.
\newcommand{\importpaperaux}[2]{%
  \IfFileExists{../../build/paper-aux/#2.aux}{%
    \externaldocument[#1]{../../build/paper-aux/#2}%
  }{%
    \IfFileExists{build/paper-aux/#2.aux}{%
      \externaldocument[#1]{build/paper-aux/#2}%
    }{%
      \IfFileExists{#2.aux}{%
        \externaldocument[#1]{#2}%
      }{}%
    }%
  }%
}

% General external reference macros
\newcommand{\paperref}[2]{\ref{#1-#2}}
\newcommand{\papereqref}[2]{(\ref{#1-#2})}

% CPSV16 source (1606.00608 / MPDO-22-12-17-2.tex) external document and shortcuts
\importpaperaux{cpsv16-}{1606.00608/MPDO-22-12-17-2-tnlean}
\newcommand{\cpsvref}[1]{\paperref{cpsv16}{#1}}
\newcommand{\cpsveqref}[1]{\papereqref{cpsv16}{#1}}

% Verdict marker: \gapnote{<kind>}{<status>}. Kind is the policy
% classification (clarification, local-correction, scope-restriction,
% unfaithful, false-source, open-gap); status is open, wip (elimination
% underway), resolved, or historical. Machine-read by the site index;
% prints a verdict line.
\newcommand{\gapnote}[2]{\par\noindent\textbf{Verdict:} #1 (#2).\par}


\title{Projective Reading of Proportional MPVs in CPSV16}
\author{}
\date{2026-05-11}

\begin{document}
\maketitle
\gapnote{clarification}{resolved}

\section{Source phrase}

Cirac--P\'erez-Garc\'ia--Schuch--Verstraete state
Theorem~\texttt{thm1} for tensors whose matrix-product vectors are
``proportional to each other''
\cite[lines~1167--1170]{Cirac2016MPDO_arXiv}. The displayed source statement
does not specify whether the proportionality scalar at each fixed length must
be nonzero. The two possible readings are
\begin{align}
    \exists c_N\in\mathbb C,\quad
    V^{(N)}(A)
        &=c_NV^{(N)}(B),
    \label{eq:cpsv16_nonzero_proportionality_reading_weak_proportionality}\\
    \exists c_N\in\mathbb C\setminus\{0\},\quad
    V^{(N)}(A)
        &=c_NV^{(N)}(B).
    \label{eq:cpsv16_nonzero_proportionality_reading_projective_proportionality}
\end{align}
The weak relation in
(\ref{eq:cpsv16_nonzero_proportionality_reading_weak_proportionality}) permits
\begin{align}
    V^{(N)}(A)
        &=0\cdot V^{(N)}(B)
    \label{eq:cpsv16_nonzero_proportionality_reading_zero_scalar}
\end{align}
whenever $V^{(N)}(A)=0$. A scalar multiple with coefficient zero does not
define equality of projective points.

\section{Formal reading}

The formal predicate
\leanid{NonzeroProportionalMPV}\textsubscript{2} requires that, for every
positive length $N\geq 1$, there be a scalar $c_N\neq 0$ such that
\begin{align}
    V^{(N)}(A)
        &=c_NV^{(N)}(B).
    \label{eq:cpsv16_nonzero_proportionality_reading_formal_relation}
\end{align}
This is the projective reading: the two nonzero vectors determine the same
point in projective space. It is also the reading used in the proof of
Theorem~\texttt{thm1} of Ref.~\cite{Cirac2016MPDO_arXiv}. The predicate
excludes the empty word because it is not a physical chain and its coefficient
records only the bond-space dimension.

The block-selection argument assumes that the overlaps of one fixed block with
all blocks on the other side tend to zero. It expands both BNT decompositions
and invokes Corollary~\texttt{Lem1} of
Ref.~\cite{Cirac2016MPDO_arXiv} to contradict proportionality. To display the
nonzero scalar used in that contradiction, set
$a_j(N)=V^{(N)}(A_j)$ and $b_\ell(N)=V^{(N)}(B_\ell)$. Let $g_A$ and $g_B$ be
the numbers of sectors, and let $\mu_j,\nu_\ell\in\mathbb C$ be their weights
in the one-representative form used by the argument. The proportionality
relation is
\begin{align}
    \sum_{j=1}^{g_A}\mu_j^N a_j(N)
        &=c_N\sum_{\ell=1}^{g_B}\nu_\ell^N b_\ell(N).
    \label{eq:cpsv16_nonzero_proportionality_reading_bnt_expansion}
\end{align}
Projecting
(\ref{eq:cpsv16_nonzero_proportionality_reading_bnt_expansion}) onto a selected
$B$-sector $k$ gives
\begin{align}
    c_N\nu_k^N\braket{b_k(N)}{b_k(N)}
        &=0.
    \label{eq:cpsv16_nonzero_proportionality_reading_projected_sector}
\end{align}
This contradicts
\begin{align}
    c_N&\neq 0,
    \label{eq:cpsv16_nonzero_proportionality_reading_nonzero_scalar}\\
    \nu_k&\neq 0,
    \label{eq:cpsv16_nonzero_proportionality_reading_nonzero_weight}\\
    \braket{b_k(N)}{b_k(N)}&\neq 0.
    \label{eq:cpsv16_nonzero_proportionality_reading_nonzero_norm}
\end{align}
Under the canonical-form hypotheses used in the source argument, the selected
BNT block states have nonzero norm on the tail of lengths used for the
contradiction. On that tail, even the weak scalar-multiple relation forces
$c_N\neq 0$. The formal predicate therefore records the projective relation
needed at the use site; it does not add an independent asymptotic assumption.

\section{Effect on the theorem}

The nonzero requirement does not change the intended conclusion of the source
theorem. It records the conventional meaning of proportional nonzero vectors.
No named predicate retains the weaker relation that permits $c_N=0$, because
(\ref{eq:cpsv16_nonzero_proportionality_reading_zero_scalar}) is not a
projective statement.

One eventual-proportionality hypothesis remains outside the named predicate.
The overlap-convergence theorem
\leanid{mpvOverlap_norm_tendsto_one_of_eventually_proportionalMPV}%
\textsubscript{2} in
\leanid{TNLean/MPS/FundamentalTheorem/Proportional.lean} assumes only that,
eventually in $N$, some scalar $c_N$ satisfies
(\ref{eq:cpsv16_nonzero_proportionality_reading_formal_relation}); it does not
assume $c_N\neq 0$. Its self-overlap convergence hypotheses imply the eventual
nonvanishing of $c_N$.

\section{Elimination plan}

No additional theorem is needed for the projective reading of
CPSV16 Theorem~\texttt{thm1}
\cite[lines~1167--1170]{Cirac2016MPDO_arXiv}. The vanishing-scalar case is a
degenerate reading outside the source convention and receives no separate
predicate.

\bibliographystyle{alpha}
\bibliography{references}

\end{document}
